\chapter{Refactorización del backend y arquitectura del módulo de capacitaciones}
\begingroup
\color{black}
\hypersetup{linkcolor=black,urlcolor=black,citecolor=black}
\lstset{keywordstyle=\color{black},commentstyle=\color{black},stringstyle=\color{black},numberstyle=\tiny\color{black}}

\section{Responsabilidad y referencia verificable}
La contribución de Adel Alejandro Merma Ccarhuarupay corresponde a arquitectura de software y desarrollo backend, principalmente HU-04 y soporte de HU-03. La implementación y revisión técnica fueron asistidas por Codex; no equivalen a revisión de otro integrante humano. Fecha de trabajo: 17 de noviembre de 2026, zona America/Lima.

La segunda etapa parte de una aplicación funcional conservada en el commit \texttt{eb00464}, mensaje \textit{feat: subir sistema funcional de capacitaciones Paz Sur}. No debe confundirse con \texttt{cf6e720}, referencia documental anterior a la primera construcción. La rama \path{feature/adel-refactor-backend} se creó desde \texttt{main}. Se conservaron los cambios locales WAL/SHM y las ramas de compañeros; las pruebas no abrieron la base de trabajo.

\section{Descripción del backend inicial}
La primera versión ya dispone de Express, rutas que comprueban permisos, servicios de negocio, Zod y SQLite mediante \texttt{node:sqlite}. Las sesiones utilizan cookie HttpOnly y las mutaciones requieren CSRF. La capacitación y su auditoría ya se guardaban dentro de una misma transacción. Estos elementos son antecedentes conservados, no logros nuevos atribuidos a esta refactorización.

\section{Problemas identificados y objetivos}
\begin{longtable}{P{0.44\textwidth}P{0.49\textwidth}}
\toprule
Problema observado & Cambio y finalidad \\
\midrule
\endhead
El servicio reúne esquemas Zod, consultas SQL, alcance de lectura, transiciones y coordinación transaccional. & Extraer validadores y repositorio para aislar responsabilidades concretas. \\
El SQL inicial sólo compara fechas y horas como texto. Acepta escrituras directas como un día inexistente o una hora fuera de rango. & Añadir restricciones temporales mediante una migración incremental, preservando filas e identificadores. \\
La regresión inicial no cubre la matriz completa de estados ni el fallo de auditoría al editar. & Comprobar contrato, atomicidad, estados y persistencia con datos ficticios aislados. \\
\bottomrule
\end{longtable}

El objetivo es mejorar la mantenibilidad y la integridad dentro de HU-03/HU-04, sin incorporar inscripciones, certificados ni otros módulos futuros. No se añaden dependencias.

\section{Principios de diseño aplicados}
Se aplica responsabilidad única: los validadores describen entradas; el repositorio contiene SQL parametrizado; el servicio coordina el caso de uso y sus reglas. Se mejora la cohesión y se mantiene bajo acoplamiento con HTTP: el servicio no depende de Express y el frontend continúa dependiendo del contrato JSON.

La conexión de base recibida como argumento ya permitía pruebas aisladas y se conserva. No se declara el cumplimiento de todos los principios SOLID ni se introducen interfaces o clases vacías para aparentarlo.

\section{Refactorización de rutas y servicios}
Las rutas existentes ya traducen HTTP y comprueban autorización, por lo que se mantienen sus contratos. El servicio conserva el control transaccional, el responsable derivado de la sesión y el bloqueo de ediciones terminales. La nueva separación ubica los esquemas en \path{server/validation/training-validation.js} y las consultas en \path{server/repositories/training-repository.js}.

\section{Acceso a datos e integridad temporal}
La migración \path{server/migrations/002_training_dates.sql} añade controles de calendario y horario al insertar o modificar capacitaciones. No reescribe la migración inicial ni reconstruye la tabla. Comprueba los registros existentes y aborta transaccionalmente si encuentra datos inválidos; no corrige valores históricos sin decisión explícita.

Las relaciones con usuarios, el código único sin distinguir mayúsculas y la capacidad entre 1 y 10000 ya existían. Su conservación se verifica con pruebas. La aplicación de 002 a una base propia requiere respaldo consistente y control de los escritores. Las pruebas de esta etapa utilizan bases independientes, nunca la base local del usuario.

La reversión comprobada en una base ficticia elimina ambos triggers y su registro de migración dentro de una transacción, sin borrar filas:
\begin{lstlisting}[language={},caption={Reversión de la migración 002 en una conexión controlada}]
BEGIN IMMEDIATE;
DROP TRIGGER trainings_validate_dates_insert;
DROP TRIGGER trainings_validate_dates_update;
DELETE FROM schema_migrations WHERE version='002_training_dates.sql';
COMMIT;
\end{lstlisting}
Ante un error previo al COMMIT corresponde ROLLBACK. El backend actual reaplica 002 al abrir la base si su versión está ausente; mantener una reversión requiere usar el código anterior. El procedimiento completo figura en \path{docs/adel/REFACTOR_BACKEND.md}.

\section{API REST y compatibilidad}
Se mantienen \texttt{POST}, \texttt{GET} y \texttt{PUT} bajo \path{/api/trainings}. POST devuelve 201; PUT y GET, 200. Los objetos conservan \texttt{data}; los listados, \texttt{data} y \texttt{meta}; los errores, \texttt{error.message} y campos cuando corresponde. PUT conserva al responsable y la fecha de creación originales.

Se mantienen 422 para entradas inválidas, 404 para registros inexistentes o fuera del alcance y 409 para duplicados o cambios de estado prohibidos. Los errores inesperados reciben 500 genérico. La revisión del contrato no introduce rutas nuevas ni exige modificar el formulario existente.

\section{Evaluación del patrón State}
El problema consiste en aceptar una transición autorizada y rechazar ediciones terminales. La primera versión ya utiliza el mapa \texttt{TRANSITIONS} en \path{server/constants.js}. Las funciones relevantes son \texttt{createTraining} y \texttt{updateTraining}.

Se decide conservar el mapa y las guardas: existen cinco estados y reglas pequeñas, sin comportamientos extensos distintos por estado. State con clases aumentaría delegación y archivos sin resolver una necesidad adicional. La alternativa directa facilita lectura y pruebas; su límite aparece si futuras reglas de cierre o certificación agregan comportamiento complejo o duplicado.

\begin{longtable}{P{0.28\textwidth}P{0.64\textwidth}}
\toprule
Estado actual & Destinos admitidos al editar \\
\midrule
\endhead
\texttt{planificada} & planificada, programada, cancelada. \\
\texttt{programada} & programada, planificada, en\_curso, cancelada. \\
\texttt{en\_curso} & en\_curso, finalizada, cancelada. \\
\texttt{finalizada} & Ninguno; tampoco admite modificar otros campos. \\
\texttt{cancelada} & Ninguno; tampoco admite modificar otros campos. \\
\bottomrule
\end{longtable}

Ejemplo: planificada, programada, en curso y finalizada forman una secuencia válida. Saltar de planificada a finalizada devuelve 409. El grafo, el bloqueo terminal y el horario diario son supuestos existentes pendientes de validación con Paz Sur; no se acredita aprobación del cliente.

\section{Comparación real del código}
El siguiente fragmento se obtuvo de \texttt{git show eb00464:server/services/training-service.js}; muestra la validación dentro del servicio original:

\begin{lstlisting}[language=JavaScript,caption={Antes: validación dentro del servicio en eb00464}]
const dateSchema = z.string().regex(/^\d{4}-\d{2}-\d{2}$/, 'Use una fecha YYYY-MM-DD.').refine(value => {
  const date = new Date(`${value}T12:00:00Z`);
  return Number.isFinite(date.getTime()) && date.toISOString().slice(0, 10) === value;
}, 'La fecha no existe.');
\end{lstlisting}

Después, la validación se consulta directamente en su fuente real:
\fragmentocodigo{server/validation/training-validation.js}{1}{40}{Después: validación separada de capacitaciones}

Antes, la persistencia estaba dentro de \texttt{updateTraining}:
\begin{lstlisting}[language=JavaScript,caption={Antes: actualización y auditoría dentro del servicio}]
db.prepare('UPDATE trainings SET code=?,title=?,description=?,start_date=?,end_date=?,start_time=?,end_time=?,modality=?,location=?,capacity=?,status=?,instructor_id=?,updated_at=? WHERE id=?').run(value.code, value.title, value.description, value.start_date, value.end_date, value.start_time, value.end_time, value.modality, value.location, value.capacity, value.status, value.instructor_id, now(), id);
audit(db, user.id, 'update', 'training', id);
return getTraining(db, id, user);
\end{lstlisting}

La transacción y la auditoría permanecen coordinadas por el servicio refactorizado; el SQL se consulta en el repositorio real:
\fragmentocodigo{server/services/training-service.js}{1}{100}{Después: coordinación del caso de uso}
\fragmentocodigo{server/repositories/training-repository.js}{1}{60}{Después: consultas parametrizadas separadas}

En el esquema anterior, \texttt{end\_date >= start\_date} y \texttt{end\_time > start\_time} sólo comparaban textos. La migración nueva refuerza esos controles sin cambiar el motor. La evidencia completa y los comandos de comparación se conservan en \path{docs/adel/REFACTOR_BACKEND.md}. El historial permite comparar versiones reales; no demuestra quién tecleó cada línea manualmente.

\section{Pruebas ejecutadas y resultados}
Antes de la refactorización se ejecutó \texttt{npm.cmd test} el 17/11/2026: 28 pruebas aprobadas y ningún fallo. Después, \texttt{npm.cmd run check} aprobó 68 pruebas, incluidos los 25 subcasos de estados, y compiló TypeScript y Vite con código cero. La advertencia de tamaño del paquete no bloqueó la compilación; las fuentes frontend y las dependencias no cambiaron en esta etapa.

La ejecución focal de migración y persistencia aprobó ocho pruebas; la revisión técnica independiente asistida por IA no identificó hallazgos bloqueantes y ejecutó 48 pruebas focales de capacitaciones y seguridad. Son resultados técnicos, no revisión de un estudiante humano. El registro completo está en \path{docs/PRUEBAS.md}.

Playwright completó cuatro casos en el primer intento, pero el runner quedó trabado y se interrumpió con código 1. La repetición con permisos ampliados finalizó correctamente: cuatro aprobadas en 8,7 segundos, código cero. Las capturas nuevas se conservan fuera del seguimiento en \path{.verification/adel-ui/}; las históricas versionadas se restauraron desde \texttt{eb00464}.

Los casos añadidos abarcan campos obligatorios, calendario y horarios, registros inexistentes, códigos duplicados, transiciones válidas e inválidas, reversión por fallo de auditoría y recuperación tras editar y reabrir SQLite. Las pruebas de migración verifican conservación de datos, rechazo de registros inválidos y aplicación incremental. No se utilizan datos personales reales.

No se encontró compilador LaTeX en PATH ni en las ubicaciones MiKTeX consultadas. Se conservan fuentes para Overleaf y su preparador. El archivo \path{docs/latex/main.pdf} ya existía en \texttt{eb00464}; no se generó un PDF actualizado con este capítulo y no se presenta el anterior como resultado de esta etapa.

El preparador Overleaf se ejecutó con código cero. Los commits reales de código y pruebas son \texttt{36b2f2d}, \texttt{1440c84} y \texttt{541a3bf}. El commit \texttt{63cdbd9} retiró sólo del índice el entorno y las bases que ya estaban rastreados, conservando sus archivos locales; no saneó ni reescribió el historial anterior.

\section{Conclusiones individuales y pendientes}
La contribución separa responsabilidades concretas y refuerza una invariante temporal demostrable, conservando el contrato del frontend y la atomicidad del caso de uso. State se evalúa y se descarta por falta de complejidad que lo justifique en esta entrega.

La implementación asistida por Codex constituye evidencia técnica, no una revisión humana ni una integración. Según la definición de terminado, siguen pendientes la revisión por otro integrante, la integración autorizada y la validación de supuestos comerciales. Los textos para Jira y el guion oral están en \path{docs/adel/}; no se fabrican fechas de Sprint, reuniones o revisiones.

\endgroup
